% ============================================================
% Section 170: 椭球状刚性势阱中的基态能量修正
% ============================================================
\section{量子力学：椭球状刚性势阱中的基态能量修正}
\label{sec:ellipsoid-well-perturbation}

\begin{modelbox}[题目：椭球状刚性势阱中的基态能量修正]
\begin{enumerate}
    \item 证明在通常的定态微扰论中，若 $H=H_0+H'$ 且 $H_0\phi_0=E_0\phi_0$，则能量一级修正为 $\Delta E_0=\langle\phi_0|H'|\phi_0\rangle$。
    \item 球形核中核子处于半径 $R$ 的球对称无限深势阱中。若核发生微小形变，势阱变为椭球形（仍为无限高势壁）：
    \[
    \frac{x^2+y^2}{b^2}+\frac{z^2}{a^2}=1,
    \]
    其中 $a\approx R(1+2\beta/3)$，$b\approx R(1-\beta/3)$，$\beta\ll1$。利用适当的变量代换和第(1)问结果，近似求椭球形核相对于球形核基态能量的变化。
\end{enumerate}
\end{modelbox}

\begin{tipbox}[考点快速分析]
\begin{itemize}
    \item \textbf{一级微扰公式}：左乘 $\langle\phi_0|$，利用正交归一性直接导出
    \item \textbf{坐标变换}：将椭球边界化为球面，形变部分进入动能算符
    \item \textbf{球对称性}：基态波函数各向同性，导致一阶修正为零
    \item \textbf{体积守恒}：$a b^2\approx R^3$，保证一级微扰项的对称结构
\end{itemize}
\end{tipbox}

\begin{derivationbox}[标准解题过程]
\textbf{1. 一级能量修正公式的证明}

设微扰后 $|\psi\rangle=|\phi_0\rangle+\lambda_1|\phi_1\rangle+\cdots$，代入 $(H_0+H')|\psi\rangle=(E_0+\Delta E_0)|\psi\rangle$。保留到一阶：
\[
H_0|\phi_0\rangle+H'|\phi_0\rangle+H_0\lambda_1|\phi_1\rangle=E_0|\phi_0\rangle+E_0\lambda_1|\phi_1\rangle+\Delta E_0|\phi_0\rangle.
\]

左乘 $\langle\phi_0|$，利用 $H_0|\phi_0\rangle=E_0|\phi_0\rangle$ 和 $\langle\phi_0|\phi_i\rangle=0\;(i\neq0)$，得
\begin{equation}
\boxed{\Delta E_0=\langle\phi_0|H'|\phi_0\rangle.}
\end{equation}

\textbf{2. 椭球势阱基态能量修正}

\textit{坐标变换}：令 $x=\dfrac{b}{R}\xi$，$y=\dfrac{b}{R}\eta$，$z=\dfrac{a}{R}\zeta$，则椭球面变为球面 $\xi^2+\eta^2+\zeta^2=R^2$。

\textit{动能算符变换}：
\[
H=-\frac{\hbar^2}{2m}\Bigl[\frac{R^2}{b^2}\Bigl(\frac{\partial^2}{\partial\xi^2}+\frac{\partial^2}{\partial\eta^2}\Bigr)+\frac{R^2}{a^2}\frac{\partial^2}{\partial\zeta^2}\Bigr].
\]

利用 $a\approx R(1+2\beta/3)$，$b\approx R(1-\beta/3)$，展开 $(1+\varepsilon)^{-2}\approx1-2\varepsilon$：
\[
\frac{R^2}{b^2}\approx1+\frac{2\beta}{3},\quad\frac{R^2}{a^2}\approx1-\frac{4\beta}{3}.
\]

于是 $H=H_0+H'$，其中 $H_0=-\dfrac{\hbar^2}{2m}\nabla_\xi^2$（球形势阱哈密顿量），微扰
\[
H'=-\frac{\hbar^2\beta}{3m}\Bigl(\frac{\partial^2}{\partial\xi^2}+\frac{\partial^2}{\partial\eta^2}-2\frac{\partial^2}{\partial\zeta^2}\Bigr).
\]

\textit{计算一级修正}：球形无限深势阱基态波函数
\[
\phi_0(\xi,\eta,\zeta)=\sqrt{\frac{2}{R}}\frac{\sin(\pi r/R)}{r},\quad r=\sqrt{\xi^2+\eta^2+\zeta^2},
\]
为球对称函数。因此三个坐标方向的二阶导数期望值相等：
\[
\Bigl\langle\frac{\partial^2}{\partial\xi^2}\Bigr\rangle=\Bigl\langle\frac{\partial^2}{\partial\eta^2}\Bigr\rangle=\Bigl\langle\frac{\partial^2}{\partial\zeta^2}\Bigr\rangle.
\]

微扰项中括号内组合为 $(1+1-2)=0$，故
\begin{equation}
\boxed{\Delta E_0=\langle\phi_0|H'|\phi_0\rangle=0.}
\end{equation}

\textit{物理讨论}：一级修正为零是基态球对称性的直接后果。体积守恒下的微小形变在哈密顿量中产生各向异性微扰，但各向同性基态的一阶响应恰好抵消。非零能量修正从二级微扰开始出现，量级为 $O(\beta^2)$。
\end{derivationbox}

\begin{formulabox}[答案汇总]
\begin{center}
\begin{tabular}{ll}
\toprule
\textbf{结论} & \textbf{结果} \\
\midrule
一级能量修正公式 & $\Delta E_0=\langle\phi_0|H'|\phi_0\rangle$ \\
椭球势阱一级修正 & $\Delta E_0=0$ \\
非零修正量级 & $O(\beta^2)$（二级微扰） \\
\bottomrule
\end{tabular}
\end{center}
\end{formulabox}

\begin{insightbox}[物理图像]
\begin{itemize}
    \item 该结果与 Jahn-Teller 效应形成对照：对于球对称基态，体积守恒的微小形变不产生一阶能量移动；但若基态存在简并（如轨道简并），则形变通常会产生一阶分裂（Jahn-Teller 定理）。
    \item 在核物理中，形变核的基态能量修正对研究核形状相变（球形 $	o$ 椭球形）至关重要，但需用到 beyond-mean-field 方法或集体模型。
    \item 本题展示了微扰论中对称性的威力：不需要具体计算积分，仅由对称性分析即可确定一级修正为零。
\end{itemize}
\end{insightbox}
